\documentclass[a4paper,12pt]{article}

% --- Pacchetti Fondamentali ---
\usepackage[utf8]{inputenc} % Codifica caratteri
\usepackage[T1]{fontenc}    % Codifica font
\usepackage[italian]{babel} % Lingua italiana
\usepackage{geometry}       % Gestione margini
\usepackage{parskip}        % Gestione spazi tra paragrafi
\usepackage{hyperref}
\geometry{a4paper, margin=2.5cm}

% --- Pacchetti Matematica e Fisica ---
\usepackage{amsmath, amssymb} % Simboli matematici
\usepackage{siunitx}          % Gestione unità di misura e numeri (FONDAMENTALE)
\sisetup{
    output-decimal-marker = {.}, % Punto come separatore decimale
    separate-uncertainty = true, % Scrive errori come (Valore +/- Errore)
    per-mode = power,            % Usa / per le unità (m/s)
    inter-unit-product = \ensuremath{\cdot} % Per scrivere le unità con il punto (N·m)
}

% --- Pacchetti Immagini e Tabelle ---
\usepackage{graphicx} % Per inserire PNG/JPG
\usepackage{float}    % Per forzare la posizione delle immagini (H)
\usepackage{booktabs} % Per tabelle professionali (\toprule, \midrule)
\usepackage{caption}  % Per personalizzare le didascalie
\usepackage{subcaption}
\usepackage[font=small, labelfont=bf]{caption}
\usepackage[font=small, labelfont=bf]{subcaption}


% --- Dati Intestazione ---
\title{\textbf{Studio rivelatori di raggi cosmici}}
\author{Filippo Audisio, Cataldo Insalaco, Telemaco Pezzoni}
\date{\today}

\begin{document}

\maketitle

% -------------------------------------------------------------------

%   F: su questa non ho niente da commentare, c'è solo da scrivere meglio qualche g^-1 ome g^{-1} e qualche cm^-1 come cm^{-1}, ma solo per bellezza

\section{Obiettivo dell'esperienza}
L'obiettivo dell'esperienza è studiare il flusso dei muoni cosmici a terra utilizzando scintillatori plastici. In particolare si vuole:
\begin{itemize}
    \item Costruire la curva della frequenza dei muoni cosmici in funzione dell'angolo zenitale, verificando la legge $f(\theta) = f_v \cdot \cos^2(\theta) + b$.
    \item Eseguire la stessa verifica tramite fit lineare.
    \item Valutare l'accordanza con il valore previsto di flusso verticale.
    \item Stimare l'efficienza di una lastra del rivelatore.
\end{itemize}

% -------------------------------------------------------------------
\section{Materiali e Metodi}

\subsection{Dotazione sperimentale}
\begin{itemize}
    \item Due lastre di scintillatore plastico $(15 \times 15) \text{ cm}$ accoppiate a SiPM.
    \item Modulo per la misura dei conteggi di coincidenza.
    \item Goniometro e struttura di supporto per i rivelatori.
    \item Chiavi a brugola per modificare l'angolo di puntamento del telescopio.
\end{itemize}

\subsection{Procedura sperimentale}
Per prima cosa sono stati collegati gli scintillatori plastici al modulo di conteggio. In seguito sono state effettuate acquisizioni di $\sim 5 \unit{min}$ per ciascun angolo zenitale, esplorando l'intervallo da -80° a +80° con passi di 10°; ad ogni fase si è segnato il numero di conteggi di coincidenza.
Infine, laddove necessario, sono stati prolungati i tempi di presa dati al fine di incrementare il campione e ridurre l'incertezza statistica.
% -------------------------------------------------------------------
\section{Dati sperimentali e Analisi}

% -------------------------------------------------------------------

\section{Conclusioni}

\end{document}